% Generated draft: requires mathematical and source review.
\documentclass[11pt]{article}
\usepackage[margin=25mm]{geometry}
\usepackage{fontspec}
\setmainfont{Latin Modern Roman}
\usepackage{amsmath,amssymb,mathtools,unicode-math}
\setmathfont{Latin Modern Math}
\usepackage{enumitem,needspace,xurl,hyperref}
\setlength{\parindent}{0pt}
\setlength{\parskip}{5pt}
\setlength{\emergencystretch}{3em}
\newcommand{\N}{\mathbb{N}}
\newcommand{\Z}{\mathbb{Z}}
\newcommand{\Q}{\mathbb{Q}}
\newcommand{\R}{\mathbb{R}}
\newcommand{\C}{\mathbb{C}}
\newcommand{\F}{\mathbb{F}}
\newcommand{\A}{\mathbb{A}}
\newcommand{\PP}{\mathbb{P}}
\newcommand{\E}{\mathbb{E}}
\newcommand{\eps}{\varepsilon}
\newcommand{\dd}{\,\mathrm{d}}
\begin{document}
\section*{UnsolvedMath ID: 6200091}
\textbf{Source problem code:} catalogue entry
\subsection{Definitions and mathematical statement}
Let $\mathbb H^n_{\C}$ be complex hyperbolic $n$-space and $PU(n,1)$ its holomorphic isometry group, the projective unitary group of a Hermitian form of signature $(n,1)$. A lattice is a discrete subgroup $\Gamma<PU(n,1)$ with finite Haar covolume; both uniform and nonuniform lattices are included.

Under the usual complex-hyperplane convention, a reflection is a nonidentity finite-order holomorphic isometry fixing a complex hyperplane pointwise and rotating the normal complex direction. A reflection-generated lattice is a lattice generated by such elements. The source does not explicitly fix this convention, including permitted reflection orders or whether point-fixing complex symmetries are meant to be included.

The proposed complex-hyperbolic analogue of Vinberg's finiteness theorem is
\[
 \exists N\in\N\ \forall n\ge N\quad
 \nexists\Gamma<PU(n,1)\text{ a lattice generated by reflections}.
\]
Thus the target is a dimension threshold uniform over all such lattices, not finiteness of the number of lattices in one fixed dimension and not the absence of arbitrary infinite-covolume reflection groups. No numerical value of $N$ is prescribed.
\subsection{Short English statement}
Prove that complex hyperbolic reflection lattices cannot exist in arbitrarily large complex dimensions. The statement includes all finite-covolume lattices, not just cocompact examples.

The supplied summary describes the high-dimensional existence question as open and mentions $n\ge10$ as an unresolved range, without a precise theorem or reference. That number is not substituted for the conjecture's unspecified threshold $N$. The OPEN-TRIAGE label and historical 2005 source are retained with current status explicitly marked for review.
\subsection{Sources}
\begin{itemize}
\item Local UnsolvedMath registry, permanent ID 6200091, source code catalogue entry, ``Boundaries of Groups and Kleinian Groups --- Problem 91,'' proposed year 2005; no proposer is specified. Dataset curation and metadata: Ulam AI and UnsolvedMath Contributors, CC BY 4.0.
\item Kapovich, \emph{Problems on Boundaries of Groups and Kleinian Groups} (2005), Problem 91, PDF page 23, \url{https://www.math.ucdavis.edu/~kapovich/EPR/problems.pdf}; normalized pdf-text extraction, accessed 2026-07-29.
\item The source invokes Vinberg's finiteness theorem without complete bibliographic details or an exact real-hyperbolic statement. The research summary's dimension-ten comment has no separate citation. Source-specific provenance marks current status and statement-reuse terms NEEDS\_REVIEW for the public author-hosted PDF, separately from the dataset licence.
\end{itemize}
\end{document}
